\documentclass[11pt,a4paper]{article}

\usepackage[utf8]{inputenc}
\usepackage[T1]{fontenc}
\usepackage{lmodern}
\usepackage{float}
\usepackage[margin=2.2cm]{geometry}
\usepackage{amsmath,amssymb}
\usepackage{booktabs}
\usepackage{enumitem}
\usepackage{xcolor}
\usepackage{titlesec}
\usepackage{mdframed}
\usepackage{hyperref}

\hypersetup{colorlinks=true, linkcolor=blue!70!black, urlcolor=blue!70!black}

\titleformat{\section}{\large\bfseries}{\thesection.}{0.5em}{}
\titleformat{\subsection}{\normalsize\bfseries}{\thesubsection.}{0.5em}{}
\titlespacing*{\section}{0pt}{1.4ex plus 0.5ex minus .2ex}{0.7ex}

\setlength{\parskip}{0.4em}
\setlength{\parindent}{0pt}

% Caixas para dicas e para perguntas da banca.
\newmdenv[linecolor=blue!50!black, linewidth=0.6pt, backgroundcolor=blue!4,
  roundcorner=3pt, skipabove=0.6em, skipbelow=0.6em,
  innertopmargin=0.5em, innerbottommargin=0.5em]{dica}
\newmdenv[linecolor=red!55!black, linewidth=0.6pt, backgroundcolor=red!3,
  roundcorner=3pt, skipabove=0.6em, skipbelow=0.6em,
  innertopmargin=0.5em, innerbottommargin=0.5em]{pergunta}

\newcommand{\Q}[1]{\textbf{P: #1}}
\newcommand{\A}[1]{\textit{R:} #1}

\begin{document}

\begin{center}
{\LARGE \textbf{Guia de Apresentação --- Projeto 3}}\\[0.3em]
{\large VAE para Fashion-MNIST: roteiro, justificativas e banco de perguntas}\\[0.2em]
\textbf{Autora:} Gabriely Di Folco Rocha
\end{center}
\vspace{-0.4em}\hrule\vspace{0.8em}

\begin{dica}
\textbf{Como usar este guia.} A apresentação é conduzida \emph{pelo relatório}.
Cada seção abaixo corresponde a uma seção do \texttt{main.pdf}: traz (i) a
mensagem de uma frase que você deve entregar, (ii) o \emph{porquê} de cada
escolha (é o que a banca avalia) e (iii) perguntas prováveis com respostas.
Não leia o guia: use-o para ensaiar. A meta é demonstrar que você entende
\emph{por que} o VAE funciona, não apenas o que você rodou.
\end{dica}

\section{Mensagem central (abertura, 30 s)}

\begin{dica}
Abra com o fio condutor, não com o dataset. Sugestão de fala:
``Treinei dois VAEs para o Fashion-MNIST com espaços latentes de 2 e 8
dimensões. O projeto é, no fundo, sobre um \textbf{compromisso entre compressão
e fidelidade}: o espaço 2D é interpretável mas reconstrói pior; o 8D reconstrói
melhor mas precisa ser projetado para ser visto. Esse trade-off, capturado pela
decomposição do ELBO em reconstrução e KL, organiza todos os resultados.''
\end{dica}

Se você fixar \emph{uma} ideia na cabeça da banca, que seja esta: \textbf{o ELBO
tem dois termos em tensão} --- reconstrução (quer usar o latente ao máximo) e KL
(quer colar o latente no \textit{prior}) --- e cada decisão do projeto mexe
nesse equilíbrio.

\section{O que é um VAE e por que não um autoencoder comum}

\textbf{Mensagem:} o VAE é um modelo \emph{generativo} probabilístico, não um
compressor. Ele aprende uma \emph{distribuição} no espaço latente, o que permite
amostrar imagens novas.

\begin{itemize}[leftmargin=1.2em, itemsep=0.1em]
\item \textbf{Autoencoder determinístico (AE):} mapeia $\mathbf{x}\to z\to
\hat{\mathbf{x}}$ com $z$ um ponto. O espaço latente fica cheio de ``buracos'':
amostrar um $z$ qualquer tende a gerar lixo. Não é gerador.
\item \textbf{VAE:} o encoder produz uma \emph{distribuição}
$q_\phi(z\mid x)=\mathcal{N}(\boldsymbol{\mu},\boldsymbol{\sigma}^2)$; o termo KL
força essas gaussianas a se aproximarem de um \textit{prior}
$\mathcal{N}(\mathbf{0},\mathbf{I})$. Isso ``preenche'' e regulariza o espaço, de
modo que amostrar $z\sim\mathcal{N}(0,I)$ e decodificar gere imagens plausíveis.
\end{itemize}

\begin{dica}
Frase de efeito: ``O termo KL é o que transforma um autoencoder em um
\emph{modelo generativo}: sem ele, o espaço latente não teria estrutura
amostrável.'' Use a Figura da varredura do plano 2D como evidência visual de que
o espaço é contínuo.
\end{dica}

\section{A função de custo (ELBO) --- domine isto}

Esta é a parte mais provável de ser cobrada. A loss minimizada é o ELBO
negativo:
\[
\mathcal{L}(\mathbf{x}) =
\underbrace{\mathbb{E}_{q_\phi(z\mid x)}[-\log p_\theta(x\mid z)]}
_{\text{reconstrução (BCE)}}
+\;\beta\,
\underbrace{D_{\mathrm{KL}}(q_\phi(z\mid x)\,\|\,p(z))}_{\text{regularização}} .
\]

\begin{itemize}[leftmargin=1.2em, itemsep=0.15em]
\item \textbf{Reconstrução} = BCE somada sobre os 784 pixels, média por imagem.
Mede o quão bem o decoder reproduz $x$. É o termo dominante na prática.
\item \textbf{KL} = forma fechada para gaussianas diagonais:
$-\tfrac12\sum_j(1+\log\sigma_j^2-\mu_j^2-\sigma_j^2)$. Penaliza o posterior por
se afastar do \textit{prior}; é um ``orçamento de informação'' do latente.
\item \textbf{$\beta=1$} = VAE padrão (Kingma \& Welling). $\beta>1$
($\beta$-VAE) força latentes mais desembaraçados ao custo de reconstrução.
\end{itemize}

\begin{dica}
Saiba explicar \emph{por que o ELBO é um limite inferior}: $\log p(x) =
\mathrm{ELBO} + D_{\mathrm{KL}}(q_\phi(z|x)\|p_\theta(z|x))$. Como a KL é
$\ge 0$, o ELBO $\le \log p(x)$. Maximizar o ELBO empurra a evidência para cima
\emph{e} aproxima o posterior aproximado do verdadeiro. (Minimizamos o negativo.)
\end{dica}

\section{O truque de reparametrização --- por que é necessário}

\textbf{Mensagem:} sem ele, não há gradiente através da amostragem.

Amostrar $z\sim\mathcal{N}(\mu,\sigma^2)$ é uma operação estocástica: não se
retropropaga por um ``sorteio''. A reparametrização escreve
$z=\mu+\sigma\odot\epsilon$, com $\epsilon\sim\mathcal{N}(0,I)$ \emph{externo} ao
grafo. A aleatoriedade vira uma entrada; $\mu$ e $\sigma$ viram funções
determinísticas e diferenciáveis. Assim o gradiente flui para o encoder.

\begin{dica}
Detalhe de implementação que impressiona: eu parametrizo \texttt{logvar}, não
$\sigma$. Motivo: $\sigma=\exp(0{,}5\cdot\text{logvar})$ é sempre positivo sem
precisar de restrição, e o log melhora a estabilidade numérica.
\end{dica}

\section{Arquitetura e escolhas de implementação}

\textbf{Mensagem:} MLP simples e simétrica, escolhas conservadoras e
justificáveis.

\begin{itemize}[leftmargin=1.2em, itemsep=0.1em]
\item \textbf{Encoder/decoder espelhados}, ReLU nas ocultas. Encoder termina em
\emph{duas} cabeças lineares ($\mu$ e $\log\sigma^2$).
\item \textbf{Sigmoide na saída} do decoder: pixels normalizados em $[0,1]$, e a
sigmoide combina naturalmente com a BCE.
\item \textbf{BCE como reconstrução:} trata cada pixel como uma intensidade em
$[0,1]$. Alternativa seria MSE (verossimilhança gaussiana); BCE costuma dar
bordas mais nítidas em imagens tipo MNIST.
\end{itemize}

\begin{pergunta}
\Q{Por que MLP e não CNN?} \A{Escopo do projeto é a análise do espaço latente,
não o estado da arte em reconstrução. A MLP é suficiente, mais leve para CPU e
mantém o foco no trade-off reconstrução/KL. Uma CNN melhoraria a reconstrução,
mas não muda a história central.}
\end{pergunta}

\section{Metodologia experimental (o que a banca mais respeita)}

\textbf{Mensagem:} separação rigorosa de dados e seleção honesta de modelo.

\begin{itemize}[leftmargin=1.2em, itemsep=0.1em]
\item \textbf{Partição:} 54\,000 treino / 6\,000 validação (semente 42); os
10\,000 de teste ficam \emph{intocados} até a avaliação final e as figuras.
\item \textbf{Seleção de topologia pela loss de validação}, nunca pelo teste.
Isso evita escolher arquitetura olhando o conjunto de avaliação.
\item \textbf{Busca em dois estágios:} candidatos treinados por até 8 épocas em
subamostra de 20\,000 (barato); o vencedor re-treinado por até 20 épocas no
treino completo. \textit{Early stopping} com \textit{patience} 3 (busca) e 6
(final).
\item \textbf{Espaço de busca:} ocultas $(256)$, $(512,256)$, $(512,256,128)$,
para cada dimensão latente (2 e 8). Adam, $\eta=10^{-3}$, lote 256.
\end{itemize}

\begin{dica}
Antecipe a crítica do atalho: ``Usei subamostra e poucas épocas \emph{apenas na
busca} para caber em CPU; o modelo final foi treinado no conjunto completo. É um
compromisso consciente entre custo e rigor, e a seleção continua baseada na
validação completa.''
\end{dica}

\section{Leitura dos resultados numéricos}

Tenha estes números na ponta da língua (por imagem):

\begin{table}[H]
\centering\small
\begin{tabular}{llrrr}
\toprule
\textbf{Latente} & \textbf{Topologia} & \textbf{Loss teste} &
\textbf{Recon.} & \textbf{KL} \\
\midrule
2D & $(512,256,128)$ & 258,2 & 251,5 & 6,7 \\
8D & $(512,256)$     & 240,7 & 227,9 & 12,7 \\
\bottomrule
\end{tabular}
\end{table}

\begin{itemize}[leftmargin=1.2em, itemsep=0.1em]
\item \textbf{8D reconstrói melhor} (227,9 vs 251,5): mais dimensões carregam
mais informação por imagem.
\item \textbf{8D tem KL maior} (12,7 vs 6,7): justamente porque usa mais o
latente, afasta-se mais do \textit{prior}. \textbf{Essa é a frase-chave do
trade-off.}
\item \textbf{Teste $\approx$ validação} em ambos: sem sobreajuste severo.
\item Na busca 2D, a loss caiu monotonicamente com a profundidade; na 8D, a
topologia \emph{intermediária} venceu --- a mais profunda não ajudou dentro do
orçamento de treino.
\end{itemize}

\begin{dica}
Se perguntarem ``esse 258 é bom?'': o valor absoluto do ELBO não tem unidade
interpretável isolada; o que importa é a \emph{comparação relativa} entre
configurações e a decomposição recon/KL. Não prometa que é o melhor possível.
\end{dica}

\section{Exploração do espaço latente e as três perguntas}

\textbf{Clusters?} Sim, no 2D: \textit{Trouser}, \textit{Bag} e o grupo de
calçados formam regiões contíguas e distintas --- \emph{de forma não
supervisionada}, pois os rótulos não entram no treino (só colorem o gráfico).

\textbf{Separação de rótulos?} Parcial: boa entre grupos visualmente distintos;
as peças de tronco (\textit{T-shirt}, \textit{Pullover}, \textit{Coat},
\textit{Shirt}, \textit{Dress}) se sobrepõem por semelhança de forma.

\textbf{Quanta variância a projeção ilustra?} No 2D, 100\% (as duas dimensões
são mostradas). No 8D, a PCA em 2 componentes cobre só \textbf{34,5\%}
(18,0\% + 16,6\%) --- dois terços da estrutura ficam nas 6 dimensões
descartadas.

\begin{dica}
O argumento mais fino do projeto: a figura 8D parece \emph{mais misturada} que a
2D, mas isso é em boa parte \textbf{artefato da projeção}, não ausência de
estrutura. Diga explicitamente: ``não concluo que o 8D separa pior; concluo que
\emph{a projeção} esconde a separação, porque só vejo 34,5\% da variância.''
Isso mostra maturidade analítica.
\end{dica}

\section{Banco de perguntas da banca (ensaie em voz alta)}

\begin{pergunta}
\Q{Por que o KL do 8D é maior que o do 2D?} \A{Mais dimensões latentes
codificam mais informação por imagem; o posterior se afasta mais do prior
$\mathcal{N}(0,I)$, elevando a KL. É o preço da melhor reconstrução.}
\end{pergunta}

\begin{pergunta}
\Q{O que é \textit{posterior collapse}? Você observou?} \A{É quando o KL vai a
zero e o latente é ignorado (o decoder gera a média). Aqui não houve: o KL
estabiliza em patamar positivo (6,7 e 12,7), indicando que o latente é usado.}
\end{pergunta}

\begin{pergunta}
\Q{Por que BCE e não MSE na reconstrução?} \A{Com pixels em $[0,1]$ e sigmoide
na saída, a BCE corresponde a uma verossimilhança de Bernoulli por pixel e tende
a produzir bordas mais nítidas. MSE (gaussiana) tenderia a borrar.}
\end{pergunta}

\begin{pergunta}
\Q{Por que usar $\mu$ (e não amostrar $z$) ao plotar o espaço latente?} \A{Na
visualização quero uma representação determinística e estável de cada imagem; a
média do posterior é o ponto mais provável. Amostrar adicionaria ruído sem valor
analítico.}
\end{pergunta}

\begin{pergunta}
\Q{Por que PCA e não t-SNE/UMAP para projetar o 8D?} \A{PCA é linear e me
permite \emph{quantificar} a variância retida (34,5\%), que é exatamente a
terceira pergunta do enunciado. t-SNE/UMAP dariam clusters mais bonitos, mas sem
uma medida honesta de quanto da estrutura estou mostrando, e podem criar
agrupamentos artificiais.}
\end{pergunta}

\begin{pergunta}
\Q{O que $\beta$ controla? E se aumentasse?} \A{$\beta$ pondera o KL. Aumentar
($\beta$-VAE) força um latente mais desembaraçado e mais colado ao prior, ao
custo de pior reconstrução --- deslocaria o trade-off para o lado da compressão.}
\end{pergunta}

\begin{pergunta}
\Q{Como você geraria imagens novas?} \A{Amostro $z\sim\mathcal{N}(0,I)$ e aplico
o decoder. A varredura do plano 2D é essencialmente isso de forma estruturada, e
mostra que o espaço é contínuo e amostrável.}
\end{pergunta}

\begin{pergunta}
\Q{Por que a loss de validação, e não a acurácia, guia a seleção?} \A{VAE é não
supervisionado; não há rótulo para prever. O critério natural é a própria
verossimilhança aproximada (ELBO) no conjunto de validação.}
\end{pergunta}

\begin{pergunta}
\Q{A topologia mais profunda nem sempre venceu. Por quê?} \A{No 8D, a
$(512,256,128)$ perdeu para a $(512,256)$. Mais capacidade não garante menor loss
dentro de um orçamento fixo de épocas/dados; pode treinar mais devagar ou
sobreajustar. A seleção por validação captura isso.}
\end{pergunta}

\begin{pergunta}
\Q{Qual a maior limitação do trabalho?} \A{Orçamento de treino reduzido (CPU,
subamostra na busca) e MLP em vez de CNN. Nada disso compromete a análise do
trade-off, mas melhores reconstruções exigiriam convolução e mais épocas.}
\end{pergunta}

\section{Checklist final (1 min antes de começar)}

\begin{itemize}[leftmargin=1.2em, itemsep=0.05em]
\item Sei desenhar o ELBO e nomear seus dois termos.
\item Sei explicar reparametrização em uma frase (gradiente pela amostragem).
\item Tenho os 4 números: recon/KL de 2D (251,5 / 6,7) e 8D (227,9 / 12,7).
\item Sei dizer por que KL$_{8D}>$KL$_{2D}$ (o trade-off).
\item Sei defender PCA pela \emph{variância quantificável} (34,5\%).
\item Sei admitir limitações (CPU, MLP) sem fragilizar a conclusão.
\end{itemize}

\end{document}
